\section{Introduction}\label{sec:intro}
\emph{Date: YYYY-MM-DD \quad Reading: Chapter 1 of \cite{rawlings2017}}

\subsection*{In plain words}
Write the main idea of the lecture in two or three sentences, as if
explaining it to a colleague.

\subsection{Key definitions}
\begin{definition}[Discrete-time system]
A system $\xkp = f(\xk,\uk)$ with state $\xk\in\R^n$ and input $\uk\in\R^m$.
\end{definition}

\begin{theorem}[Example statement]\label{thm:example}
For the linear system $\xkp = A\xk + B\uk$, the cost
$J = \sum_{k=0}^{N-1}\bigl(\xk\transp Q\xk + \uk\transp R\uk\bigr)$
is minimized by a linear feedback law.
\end{theorem}
\begin{proof}
Sketch the proof here, then fill in the details.
\end{proof}

\begin{remark}
Notes you want to remember go in remarks. See also \cite{mayne2000}.
\end{remark}

\subsection{Questions and ideas}
\begin{itemize}
  \item \todo{Question that remained after the lecture.}
  \item Idea: could assumption X be relaxed?
\end{itemize}
